\chapter{Introduction}

\section{Background}
Food waste represents one of the most critical socio-economic, ecological, and ethical paradoxes of our time. According to the United Nations Environment Programme (UNEP) Food Waste Index Report (2024), more than 1.05 billion tonnes of edible food are wasted globally each year, amounting to nearly one-fifth of all food available to consumers. Crucially, domestic households account for approximately 60\% of this total wastage—surpassing the collective losses generated by retail, hospitality, and food service sectors combined.

At the household level, food waste is rarely the result of intentional disposal. Instead, it stems from systemic culinary friction, disorganization, and decision paralysis. Home cooks regularly stock fresh produce, dairy, and perishable condiments in refrigerators and cupboards. Over the course of busy weeks, visual occlusion (items pushed behind containers), lack of structured inventory visibility, and uncertainty regarding recipe formulation lead to edible ingredients spoiling unconsumed. Simultaneously, households face escalating food inflation and environmental degradation, as decomposed organic waste in municipal landfills releases potent methane emissions accelerating global climate change.

Recent breakthroughs in multi-modal artificial intelligence (AI), computer vision, and large language models (LLMs) provide an unprecedented opportunity to address this domestic crisis. By equipping consumer web browsers with neural visual perception and culinary reasoning, software systems can automate the cognitive overhead of inventory tracking and empower users to transform seemingly random, leftover ingredients into wholesome gourmet meals.

\section{Motivation}
The inspiration behind developing \textbf{PlateWise} stems from the ubiquitous modern dilemma: \textit{“A refrigerator filled with items, yet nothing to eat.”} Traditional recipe search engines and cookbook applications fail to solve this problem because they require the cook to begin with a predetermined dish title (e.g., “Lasagna” or “Butter Chicken”) and subsequently purchase a strict list of ingredients. When home cooks attempt to search by available items on existing recipe websites, they are forced to manually type dozens of individual ingredient names through tedious search bars and checkboxes—an experience with high friction that users quickly abandon.

Furthermore, conventional recipe aggregators assume ideal pantry conditions. If a recipe calls for cream cheese, heavy cream, or specialty herbs that the user lacks, the user is left stranded without intelligent alternatives, ultimately ordering takeout while their perishable vegetables spoil in the crisper drawer.

The motivation of this project is driven by the imperative to:
\begin{enumerate}
    \item \textbf{Eliminate Manual Entry Friction:} Enable zero-effort pantry capture by allowing users to simply take a photograph or drag-and-drop an image of their refrigerator shelves, pantry, or grocery haul.
    \item \textbf{Democratize High-Throughput Vision AI:} Leverage state-of-the-art vision models on ultra-fast inference hardware to accurately identify an entire array of diverse produce, sauces, dairy, and spices simultaneously in real time.
    \item \textbf{Deliver Context-Aware Culinary Reasoning:} Provide personalized recipe suggestions that dynamically adapt to the user's specific context—such as time of day, current weather season, and regional cuisine style.
    \item \textbf{Bridge Ingredient Gaps with Algorithmic Auditing:} Formulate a mathematical audit score that highlights missing components and provides scientifically sound, practical culinary substitutions.
    \item \textbf{Close the Kitchen Action Loop:} Guide the user through active preparation with synchronized video instruction and interactive step-by-step checklists, permanently logging prepared recipes in an offline-resilient digital vault.
\end{enumerate}

\section{Problem Statement}
Existing kitchen utility platforms and recipe search websites lack automated visual perception, contextual adaptability, and intelligent ingredient auditing. Specifically, current systems suffer from the following fundamental limitations:
\begin{enumerate}
    \item They necessitate exhaustive manual text input of household groceries, introducing cognitive fatigue and high attrition rates.
    \item Traditional computer vision tools in food recognition are restricted to single-item classification (identifying only one isolated fruit or vegetable per scan) and catastrophically fail when presented with realistic, multi-item refrigerator shelves.
    \item Standard recipe repositories present static, rigid recipe cards that cannot dynamically adapt to missing ingredients or recommend realistic kitchen substitutions.
    \item Applications rely on heavy, expensive cloud database tiers that introduce latency, require intrusive account creation, and lose operational state during serverless cold starts.
\end{enumerate}

Therefore, the problem addressed in this project is articulated as follows:
\begin{quote}
\textit{“To design, develop, and deploy an individual, cloud-native, full-stack AI web platform capable of autonomously parsing multi-ingredient visual imagery from consumer refrigerators, applying context-aware generative reasoning to produce personalized recipes, mathematically auditing missing ingredients with smart substitutions, and providing an interactive culinary workstation with zero financial infrastructure overhead.”}
\end{quote}

\section{Objectives}
The core objectives of the PlateWise project are:
\begin{enumerate}
    \item \textbf{Multi-Modal Visual Recognition:} Develop an asynchronous image ingestion pipeline supporting drag-and-drop, direct upload, and realistic presets capable of detecting 8 to 16+ distinct produce and grocery items from a single photograph in under 2 seconds.
    \item \textbf{Contextual Recipe Synthesis:} Integrate high-reasoning large language models to synthesize 4 structured, gourmet recipes tailored to the detected ingredients, constrained by meal category (Breakfast, Lunch, Dinner), seasonal dynamics (Monsoon, Summer, Winter), and culinary traditions (Indian Comfort, Mediterranean, Continental).
    \item \textbf{Mathematical Pantry Auditing:} Implement an algorithmic audit module that calculates an exact percentage match ($0\%$ to $100\%$) between available ingredients and recipe requirements, flagging missing items and providing practical pantry substitutes.
    \item \textbf{Interactive Kitchen Workstation:} Create an interactive cooking screen that dynamically queries and embeds relevant YouTube masterclass cooking tutorials alongside an interactive, strike-through preparation checklist.
    \item \textbf{Offline-Resilient Dual-Tier Storage:} Architect a persistent storage layer combining Vercel serverless edge functions with browser-local \texttt{localStorage} mirroring to prevent data loss across cold starts.
    \item \textbf{Public Cloud Deployment:} Deliver a production-grade, responsive web application hosted on Vercel with automated continuous deployment from GitHub.
\end{enumerate}

\section{Scope of the Project}
PlateWise is engineered as a modern, accessible web application optimized for desktop, tablet, and mobile devices. 

\subsection*{Functional Scope}
\begin{itemize}
    \item Neural image capture and client-side downscaling of refrigerator and pantry photos.
    \item Multi-item visual detection via Groq Cloud's vision LPUs with zero-failure candidate fallbacks.
    \item Active pantry management with instant add, delete, clear, and category filtering across 70+ household staples.
    \item Parametric recipe generation across 3 meal times, 3 seasonal profiles, and customizable cuisine vectors.
    \item Deep ingredient auditing with missing item identification and replacement recommendations.
    \item Video tutorial playback and step tracking during active cooking.
    \item Historical meal vault archiving with “Cook Again” replay functionality.
\end{itemize}

\subsection*{Technical \& Economic Scope}
The platform is designed to operate strictly within zero-cost hobby and free tiers (Vercel Hobby Tier and Groq Free Developer Tier), demonstrating that enterprise-grade AI applications can be engineered sustainably without enterprise capital expenditure.

\section{Proposed Solution}
The proposed PlateWise architecture decomposes the culinary journey into five coordinated stages:
\begin{enumerate}
    \item \textbf{Stage 01 -- Visual Scan:} Client-side canvas compression downscales high-resolution camera photos to a standardized 1024px boundary, transmitting base64 JSON payloads to \texttt{/api/vision-scan}. The vision engine classifies all visible produce and groceries into an active inventory.
    \item \textbf{Stage 02 -- Contextual Suggestion:} The user selects temporal, seasonal, and regional parameters. The serverless suggestion endpoint (\texttt{/api/recommendations}) generates 4 customized recipes utilizing only available or easily substitutable ingredients.
    \item \textbf{Stage 03 -- Pantry Audit:} The selected dish is cross-referenced with the user's pantry through \texttt{/api/audit-ingredients}, quantifying compatibility and detailing exact ingredient substitutions.
    \item \textbf{Stage 04 -- Interactive Cooking:} An intuitive cooking station dynamically binds YouTube tutorial video streams and provides interactive step progression.
    \item \textbf{Stage 05 -- Culinary Vault:} Finished meals are logged into the persistent vault, closing the zero-waste feedback loop.
\end{enumerate}

\section{Methodology Overview}
The development methodology followed an iterative, test-driven approach:
\begin{enumerate}
    \item \textbf{Empirical Vision Probing:} Benchmarking visual language models on complex, cluttered refrigerator scenes to establish candidate fallback chains (\texttt{qwen/qwen3.8-27b} $\rightarrow$ \texttt{llama-4-scout} $\rightarrow$ \texttt{llama-3.2-11b-vision}).
    \item \textbf{Serverless Edge Architecture:} Implementing stateless Vercel Serverless Functions in Node.js with rigorous CORS, MIME type validation, and error boundaries.
    \item \textbf{Frontend Component Engineering:} Building a responsive single-page application using React 19, Vite, Framer Motion, and Lenis smooth momentum scrolling.
    \item \textbf{Dual-Storage Synchronization:} Hardening client state with LocalStorage mirroring to maintain user history across serverless cold starts.
    \item \textbf{Automated Continuous Deployment:} Configuring Git-integrated Vercel deployment pipelines for deterministic production builds.
\end{enumerate}

\section{Project Report Organization}
The remainder of this report is structured as follows:
\begin{itemize}
    \item \textbf{Chapter 2 (Literature Review)} examines existing smart kitchen platforms, computer vision in gastronomy, and generative LLMs for recipe formulation.
    \item \textbf{Chapter 3 (Problem Definition and Scope)} provides comprehensive functional, non-functional, and user persona specifications.
    \item \textbf{Chapter 4 (System Requirement Specification)} details system block diagrams, UML use case specifications, and error validation mechanisms.
    \item \textbf{Chapter 5 (Proposed Methodology and Implementation)} outlines the end-to-end mathematical models, vision parsing pipelines, and serverless edge handlers.
    \item \textbf{Chapter 6 (Testing, Results, and Conclusion)} presents experimental latency and accuracy benchmarks, summary benefits, and future avenues of research.
    \item \textbf{Appendices} provide originality verification, open-source license compliance, and detailed economic cost analysis.
\end{itemize}
